\documentclass[11pt]{article}
\usepackage[margin=23mm]{geometry}
\usepackage{amsmath,amssymb,amsthm}
\usepackage[T1]{fontenc}
\usepackage{lmodern}
\usepackage{tikz-cd}
\usepackage[colorlinks=true,urlcolor=blue,linkcolor=blue,hypertexnames=false]{hyperref}
\newcommand{\I}{\mathbb I}
\newcommand{\C}{\mathbf C}
\newcommand{\Nm}{\operatorname{Nm}}
\newcommand{\Gal}{\operatorname{Gal}}
\newcommand{\Res}{\operatorname{Res}}
\newcommand{\Inf}{\operatorname{Inf}}
\newcommand{\HT}[2]{H_T^{#1}(#2)}
\newtheorem*{target}{Theorem VII.5.1}
\newenvironment{authorresult}[2]{\par\medskip\noindent\textbf{#1 VII.#2.}\ \itshape}{\par\medskip}
\title{The second inequality for id\`ele classes\\The author's algebraic proof}
\author{J. S. Milne\\\small Source: \emph{Class Field Theory}, v4.03, 6 August 2020}
\date{}
\begin{document}
\maketitle
\noindent\small
Private reading transcription of Theorem VII.5.1 and the supplied algebraic
proof in VII\S\S5--6. The target and its proof occupy physical PDF pages
221--228 (printed pages 212--219); the supporting VII.4.5--4.7 below are on
physical pages 220--221. The complete original work is retained at
\url{https://www.jmilne.org/math/CourseNotes/CFT.pdf}.
Mathematical expressions and the author's argument are retained, with
connective prose and layout normalized. This is a transcription of published
PDF material, not author-native TeX or a new proof.
\normalsize

% BEGIN CONTEXT VII.5.1 NOTATION
\section*{Notation and the preceding source results}
For a number field $K$, $\I_K$ is its id\`ele group and
$\C_K=\I_K/K^\times$. The diagonal embedding of $K^\times$ is understood.
For a finite set $S$ containing the infinite primes, the author's notation is
\[
 \I_{K,S}=\prod_{v\in S}K_v^\times\times\prod_{v\notin S}U_v,
 \qquad U_v=\mathcal O_v^\times\quad(v\text{ finite}),
\]
and $\I_S$ abbreviates $\I_{K,S}$ when the base field is fixed.
$U(S)$ is the group of $S$-units. A local completion above $v$ is denoted by
$L_w$. Norm maps on id\`eles and on id\`ele classes use the same symbol.
The absolute values are those of the author's product formula; at a complex
place the normalized value is the squared usual complex absolute value, as
also seen in the calculation in Proposition VII.6.8 below.

The proof uses local class field theory, Kummer theory, the unit theorem,
weak approximation, and the basic cohomological restriction and
inflation--restriction sequences. The author's already established
Theorem VII.4.3 says that $h(\C_L)=[L:K]$ for a finite cyclic extension.
Its immediate Corollary VII.4.4, the first inequality, is
\[
 (\I_K:K^\times\Nm_{L/K}\I_L)\ge [L:K].
\]
Indeed, if $h(\C_L)=[L:K]$, its numerator is at least $[L:K]$.
The author also uses that $\Nm_{L/K}\I_L$ is open in $\I_K$ (VII.2.8).
These are preceding source results, not additional counted targets.
% END CONTEXT VII.5.1 NOTATION

% BEGIN SUPPORT VII.4.5-4.7
\begin{authorresult}{Lemma}{4.5}
Let $L$ be a finite solvable extension of $K$ (that is, a finite Galois
extension with solvable Galois group). If there is a subgroup $D$ of $\I_K$
such that
\begin{enumerate}
\item[(a)] $D\subset\Nm_{L/K}\I_L$, and
\item[(b)] $K^\times D$ is dense in $\I_K$,
\end{enumerate}
then $L=K$.
\end{authorresult}
\begin{proof}
If $L\ne K$, then there is a subfield $K'$ of $L$ that is cyclic over $K$
and different from $K$. Then
\[
 D\subset\Nm_{L/K}(\I_L)
 =\Nm_{K'/K}\bigl(\Nm_{L/K}(\I_L)\bigr)
 \subset\Nm_{K'/K}(\I_{K'}).
\]
\footnote{Source notation note: on physical page 220 (printed page 211),
the inner norm in this display is printed with the subscript $L/K$.
That published subscript is retained here as an original notation issue.}
Therefore $K^\times\Nm_{K'/K}\I_{K'}$ is dense in $\I_K$. As it is a union
of translates of $\Nm_{K'/K}\I_{K'}$, it is open by VII.2.8; being a
subgroup of $\I_K$, it is also closed. Hence
$K^\times\Nm_{K'/K}\I_{K'}=\I_K$, and the first inequality implies $K'=K$.
\end{proof}

\begin{authorresult}{Proposition}{4.6}
Let $L$ be a finite solvable extension of $K$. If $L\ne K$, there are
infinitely many primes of $K$ that do not split completely in $L$.
\end{authorresult}
\begin{proof}
Suppose there are only finitely many, and let $S\supset S_\infty$ be a finite
set of primes of $K$ including all those that do not split completely in $L$.
Apply the lemma with
\[
 D=\I^S\overset{\mathrm{def}}=
 \{(a_v)\mid a_v=1\text{ for all }v\in S\}.
\]
For $w\mid v\notin S$, $L_w=K_v$, so $D\subset\Nm(\I_L)$.
Let $\mathbf a=(a_v)\in\I_K$. By weak approximation (ANT, 7.20), there is
$b\in K^\times$ very close to $a_v$ in $K_v$ for all $v\in S$.
Choose $\mathbf a'\in\I^S$ whose $v$-component is that of
$b^{-1}\mathbf a$ for every $v\notin S$. Then $b\mathbf a'$ is close to
$\mathbf a$ in $\I_K$. Thus $K^\times D$ is dense in $\I_K$.
\end{proof}

\begin{authorresult}{Proposition}{4.7}
Let $L/K$ be a finite solvable extension with Galois group $G$, and let $T$
be a finite set of prime ideals containing those that ramify in $L$.
The Frobenius elements $(\mathfrak P,L/K)$ for $\mathfrak P\notin T$
generate $G$.
\end{authorresult}
\begin{proof}
After enlarging $T$, we may suppose it is stable under $G$. The subgroup
$H$ generated by the Frobenius elements at the $\mathfrak P\notin T$ is then
normal. Let $E=L^H$. Recall (V, 1.11) that
\[
 (\mathfrak P\cap E,E/K)=(\mathfrak P,L/K)|_E,
\]
which is the identity map. Therefore every prime $\mathfrak p$ of $K$ not
lying under a prime of $T$ splits in $E$, which shows that $E=K$.
By the main theorem of Galois theory, $H=G$.
\end{proof}
% END SUPPORT VII.4.5-4.7

% BEGIN PRIMARY STATEMENT VII.5.1
\section*{The target statement}
\begin{target}
Let $L/K$ be a Galois extension of number fields with Galois group $G$. Then
\begin{enumerate}
\item[(a)] \emph{(Second inequality)} The index
$(\I_K:K^\times\Nm(\I_L))$ is finite and divides $[L:K]$.
\item[(b)] $H^1(G,\C_L)=0$.
\item[(c)] $H^2(G,\C_L)$ is finite and its order divides $[L:K]$.
\end{enumerate}
\end{target}
% END PRIMARY STATEMENT VII.5.1

% BEGIN PRIMARY PROOF VII.5.1 REDUCTIONS
\section*{The author's proof: reductions}
\begin{authorresult}{Lemma}{5.2}
If $G$ is cyclic, statements (a), (b), and (c) are equivalent, and
\[
 (\I_K:K^\times\Nm(\I_L))=(H^2(G,\C_L):1)=[L:K].
\]
\end{authorresult}
\begin{proof}
Without restriction on $G$,
\[
 \I_K/K^\times\Nm_{L/K}(\I_L)
 \simeq\C_K/\Nm_{L/K}(\C_L)=\HT0{G,\C_L}.
\]
If $G$ is cyclic, its cohomology is periodic, so
$\HT0{G,\C_L}\simeq H^2(G,\C_L)$. This makes (a) and (c) equivalent.
Theorem VII.4.3 states that $h(\C_L)=[L:K]$, so each of (a) and (c) is
equivalent to (b).
\end{proof}

\begin{authorresult}{Lemma}{5.3}
It suffices to prove the theorem when $G$ is a $p$-group, $p$ prime.
\end{authorresult}
\begin{proof}
Recall (II, 1.33) that if $H$ is the Sylow $p$-subgroup of $G$, then for every
$G$-module $M$ the maps
\[
 \Res:\HT r{G,M}\longrightarrow\HT r{H,M}
\]
are injective on the $p$-primary components. Therefore, if the theorem holds
for $L/L^H$, then $p$ does not divide the order of $H_T^1(G,\C_L)$, and the
power of $p$ dividing the orders of $H_T^0(G,\C_L)$ and $H_T^2(G,\C_L)$ is
less than the power of $p$ dividing $[L:K]$. Applying this for every $p$ gives
the lemma.
\end{proof}

\begin{authorresult}{Lemma}{5.4}
It suffices to prove the theorem when $G$ is cyclic of prime order $p$.
\end{authorresult}
\begin{proof}
After the last lemma, assume $G$ is a $p$-group. Induct on $(G:1)$.
Because $G$ is a $p$-group, it has a normal subgroup $H$ of index $p$
(GT, 4.17). Put $K'=L^H$ and consider the exact sequence (II, 1.34)
\[
 0\longrightarrow H^1(G/H,\C_{K'})
 \xrightarrow{\Inf}H^1(G,\C_L)
 \xrightarrow{\Res}H^1(H,\C_L).
\]
By assumption $H^1(G/H,\C_{K'})=0$, and by induction $H^1(H,\C_L)=0$.
Thus $H^1(G,\C_L)=0$, proving (b). Because $H^1(H,\C_L)=0$, the sequence
\[
 0\longrightarrow H^2(G/H,\C_{K'})
 \longrightarrow H^2(G,\C_L)\longrightarrow H^2(H,\C_L)
\]
is exact, which proves (c). Finally,
\begin{align*}
 (\C_K:\Nm_{L/K}\C_L)
 &= (\C_K:\Nm_{K'/K}\C_{K'})\\
 &\quad\cdot(\Nm_{K'/K}\C_{K'}:\Nm_{L/K}\C_L).
\end{align*}
This divides $p[L:K']$ because $\Nm_{K'/K}$ defines a surjection
\[
 \C_{K'}/\Nm_{L/K'}\C_L
 \longrightarrow\Nm_{K'/K}\C_{K'}/\Nm_{L/K}\C_L.
\]
\end{proof}

To finish Theorem VII.5.1 using Lemma VII.5.2, it remains to prove the second
inequality for cyclic extensions of degree $p$. The author notes a route
through VI.4.9, and then gives the following algebraic proof independent of
Chapter VI. The algebraic route is the one retained here.
% END PRIMARY PROOF VII.5.1 REDUCTIONS

% BEGIN PRIMARY PROOF VII.5.1 ROOTS OF UNITY
\section*{The algebraic proof of the second inequality}
We prove the second inequality when $L/K$ is cyclic of prime degree $p$.
\begin{authorresult}{Lemma}{6.1}
It suffices to prove the second inequality when $K$ contains a $p$th root of $1$.
\end{authorresult}
\begin{proof}
Let $\zeta$ be a primitive $p$th root of $1$ in a fixed algebraic closure of
$K$ containing $L$. Put $K'=K[\zeta]$ and $L'=K'L=L[\zeta]$.
Then $[K':K]=m\mid p-1$, which is relatively prime to $p$, so $K'\cap L=K$:
\[
\begin{tikzcd}[row sep=large,column sep=huge]
 & L'\arrow[dl,dash,"m"']\arrow[dr,dash,"p"] & \\
 L\arrow[dr,dash,"p"'] && K'\arrow[dl,dash,"m"] \\
 & K &
\end{tikzcd}
\]
The map $\Gal(L'/K)\longrightarrow\Gal(L/K)\times\Gal(K'/K)$ is an
isomorphism. Consider the diagram:
\[
\begin{tikzcd}[column sep=large,row sep=large]
\C_L\arrow[r,"\Nm_{L/K}"]\arrow[d,"i_L"']&
\C_K\arrow[r]\arrow[d,"i_K"']&\C_K/\Nm\C_L\arrow[r]\arrow[d]&0\\
\C_{L'}\arrow[r,"\Nm_{L'/K'}"]\arrow[d,"\Nm_{L'/L}"']&
\C_{K'}\arrow[r]\arrow[d,"\Nm_{K'/K}"']&\C_{K'}/\Nm\C_{L'}\arrow[r]\arrow[d]&0\\
\C_L\arrow[r,"\Nm_{L/K}"]&\C_K\arrow[r]&\C_K/\Nm\C_L\arrow[r]&0
\end{tikzcd}
\]
Here $i_L,i_K$ are induced by $\I_L\hookrightarrow\I_{L'}$ and
$\I_K\hookrightarrow\I_{K'}$. The maps $\Nm_{L/K}$ and $\Nm_{L'/K'}$ are
\[
 x\longmapsto\sum\sigma x,
 \qquad\sigma\in\Gal(L/K)\simeq\Gal(L'/K'),
\]
and $\Nm_{L'/L}$ and $\Nm_{K'/K}$ are
\[
 x\longmapsto\sum\sigma x,
 \qquad\sigma\in\Gal(L'/L)\simeq\Gal(K'/K).
\]
The left squares commute, which implies that the rest of the diagram exists.
The composites $\Nm_{L'/L}\circ i_L$ and $\Nm_{K'/K}\circ i_K$ are both
multiplication by $m$. Hence the composite
\[
 \C_K/\Nm\C_L\longrightarrow\C_{K'}/\Nm\C_{L'}
 \longrightarrow\C_K/\Nm\C_L
\]
is also multiplication by $m$ and is therefore an isomorphism: $p$th powers
in $\C_K$ are norms, so $\C_K/\Nm\C_L$ is killed by $p$.
In particular, the second map is surjective, and
\[
 (\C_K:\Nm\C_L)\text{ divides }(\C_{K'}:\Nm\C_{L'}),
\]
which by assumption divides $p$.
\end{proof}
% END PRIMARY PROOF VII.5.1 ROOTS OF UNITY

% BEGIN PRIMARY PROOF VII.5.1 AUXILIARY EXTENSION AND PRIMES
We prove the second inequality when $K$ contains a primitive $p$th root of
$1$ and $L$ is a finite abelian extension of exponent $p$, with Galois group
$G$. Let $[L:K]=p^r$, so $G\simeq(\mathbb Z/p\mathbb Z)^r$.
By Kummer theory (the appendix to Chapter VII), there are $a_1,\ldots,a_r\in K$
such that
\[
 L=K[a_1^{1/p},\ldots,a_r^{1/p}].
\]
Choose a finite set $S$ of primes of $K$ such that
\begin{enumerate}
\item[(a)] $S$ contains the infinite primes;
\item[(b)] $S$ contains all divisors of $p$;
\item[(c)] all $a_i$ are $S$-units;
\item[(d)] $S$ contains a set of generators for the ideal class group of $K$,
so $\I_K=\I_{K,S}K^\times$ (VII.4.2).
\end{enumerate}
By A.5, (b) and (c) imply that $S$ contains all primes ramifying in $L$.
As usual, $U(S)$ consists of the elements of $K^\times$ that are units at all
primes outside $S$. The unit theorem says
\[
 U(S)\simeq\mathbb Z^{s-1}\times U(S)_{\mathrm{tors}},\qquad s=|S|.
\]
$U(S)_{\mathrm{tors}}$ is cyclic, and its order is divisible by $p$ because
it contains a primitive $p$th root of $1$. Therefore
\[
 U(S)/U(S)^p\simeq(\mathbb Z/p\mathbb Z)^s.
\]
Let $M=K[U(S)^{1/p}]$. This is the Kummer extension corresponding to
\[
 U(S)K^{\times p}/K^{\times p}
 \simeq U(S)/(U(S)\cap K^{\times p})
 =U(S)/U(S)^p\simeq(\mathbb Z/p\mathbb Z)^s.
\]
Thus there are extensions
\[
 M\overset{p^t}{\supset}L\overset{p^r}{\supset}K,\qquad r+t=s.
\]

\begin{authorresult}{Lemma}{6.2}
There is a finite set $T$ of primes of $K$, disjoint from $S$, such that
\[
 \{(\mathfrak p_v,M/K)\mid v\in T\}
\]
is a basis for $\Gal(M/L)$, regarded as an $\mathbb F_p$-vector space.
\end{authorresult}
\begin{proof}
$S$ contains all primes of $K$ ramifying in $M$ (A.5). For a prime $w$ of
$M$ above $v\notin S$, $M_w/K_v$ is therefore unramified, hence cyclic of
exponent $p$. It is either cyclic of order $p$ or trivial. Thus if
$M_w\ne L_w$, then $L_w=K_v$; here the author uses $w$ for both the prime
of $M$ and the prime of $L$ it divides.

By VII.4.7 there is a finite set $T'$ of primes of $M$, none above a prime
in $S$, such that $(\mathfrak p_w,M/L)$, $w\in T'$, generate $\Gal(M/L)$.
After replacing $T'$ by a subset, these Frobenius elements form a basis.
Let $w\in T'$ divide $w_L$ in $L$ and $w_K$ in $K$. Because
$M_w\ne L_{w_L}$, we have $L_{w_L}=K_{w_K}$, and so
$(\mathfrak p_w,M/L)=(\mathfrak p_w,M/K)$.
Thus $T=\{w_K\mid w\in T'\}$ has the required property.
\end{proof}

The order of $T$ is $t$, where $p^t=[M:L]$. If $w$ is a prime of $M$ above
$v\in T$, then $L_w=K_v$.
\begin{authorresult}{Lemma}{6.3}
With this notation, $a\in U(S)$ becomes a $p$th power in $L$ if and only if
it becomes a $p$th power in $K_v$ for every $v\in T$.
\end{authorresult}
\begin{proof}
If $a$ becomes a $p$th power in $L$, then it does so in $L_w$ for every $w$.
For $w\mid v\in T$, $L_w=K_v$, so it becomes a $p$th power in $K_v$.
Conversely, let $a\in U(S)$. Then $a^{1/p}\in M$. If $a^{1/p}\in K_v$
for all $v\in T$, it is fixed by $(\mathfrak p_v,M/K)$ for all $v\in T$,
and hence by $\Gal(M/L)$. Thus $a^{1/p}\in L$ and $a\in L^p$.
\end{proof}
% END PRIMARY PROOF VII.5.1 AUXILIARY EXTENSION AND PRIMES

% BEGIN PRIMARY PROOF VII.5.1 NORM SUBGROUP AND INDEX REDUCTION
\begin{authorresult}{Lemma}{6.4}
The subgroup
\[
 E=\prod_{v\in S}K_v^{\times p}\times
   \prod_{v\in T}K_v^{\times}\times
   \prod_{v\notin S\cup T}U_v
\]
of $\I_K$ is contained in $\Nm_{L/K}(\I_L)$.
\end{authorresult}
\begin{proof}
Let $\mathbf a=(a_v)\in E$. We show that every component $a_v$ is a norm.
For $v\in S$, local class field theory gives
\[
 K_v^\times/\Nm L_w^\times\xrightarrow{\sim}\Gal(L_w/K_v).
\]
The second group is killed by $p$, so the first is also, and
$K_v^{\times p}\subset\Nm L_w^\times$.
For $v\in T$, $L_w=K_v$, so every element of $K_v$ is a norm from $L_w$.
For $v\notin S\cup T$, $L_w/K_v$ is unramified, so the norm map
$U_w\longrightarrow U_v$ is surjective (III, 1.2).
\end{proof}

Now
\[
 (\C_K:\Nm\C_L)=(\I_K:K^\times\Nm\I_L),
\]
which divides $(\I_K:K^\times E)$. It remains to show
\[
 \boxed{(\I_K:K^\times E)\mid p^r.}
\]
But $\I_K=K^\times\I_S=K^\times\I_{S\cup T}$, and hence
\[
 (\I_K:K^\times E)=(K^\times\I_{S\cup T}:K^\times E).
\]

\begin{authorresult}{Lemma}{6.5}
Let $A,B,C$ be subgroups of an abelian group with $A\supset B$. Then
\[
 (AC:BC)(A\cap C:B\cap C)=(A:B),
\]
in the sense that if two of these indices are finite, then so is the third
and the equality holds.
\end{authorresult}
\begin{proof}
In the following commutative diagram, the columns and the first two rows
are evidently exact. It follows, for example by the snake lemma, that the
bottom row is exact, which implies the statement.
\[
\begin{tikzcd}[column sep=small,row sep=large]
 &0\arrow[d]&0\arrow[d]&0\arrow[d]&\\
0\arrow[r]&B\cap C\arrow[r]\arrow[d]&B\arrow[r]\arrow[d]&BC/C\arrow[r]\arrow[d]&0\\
0\arrow[r]&A\cap C\arrow[r]\arrow[d]&A\arrow[r]\arrow[d]&AC/C\arrow[r]\arrow[d]&0\\
0\arrow[r]&\dfrac{A\cap C}{B\cap C}\arrow[r]\arrow[d]&A/B\arrow[r]\arrow[d]&AC/BC\arrow[r]\arrow[d]&0\\
 &0&0&0&
\end{tikzcd}
\]
\end{proof}

Apply the lemma with $A=\I_{S\cup T}$, $B=E$, $C=K^\times$. Then
\[
 (\I_K:K^\times E)
 =\frac{(\I_{S\cup T}:E)}{(U(S\cup T):K^\times\cap E)}.
\]
\begin{authorresult}{Lemma}{6.6}
With the above notation, $(\I_{S\cup T}:E)=p^{2s}$.
\end{authorresult}
\begin{authorresult}{Lemma}{6.7}
With the above notation, $(U(S\cup T):K^\times\cap E)=p^{s+t}$.
\end{authorresult}
Since $r+t=s$, this will prove the boxed formula.
% END PRIMARY PROOF VII.5.1 NORM SUBGROUP AND INDEX REDUCTION

% BEGIN PRIMARY PROOF VII.5.1 LOCAL POWER INDEX
\begin{proof}[Proof of Lemma VII.6.6]
Clearly $(\I_{S\cup T}:E)=\prod_{v\in S}(K_v^\times:K_v^{\times p})$.
There are $s$ primes in $S$, and $K$ contains $p$ distinct $p$th roots of $1$.
The next proposition therefore gives
\[
 (\I_{S\cup T}:E)=\frac{p^{2s}}{\prod_{v\in S}|p|_v}.
\]
By assumption $S$ contains every prime at which $|p|_v\ne1$, so
\[
 \prod_{v\in S}|p|_v=\prod_{\text{all }v}|p|_v=1
\]
by the product formula.
\end{proof}

\begin{authorresult}{Proposition}{6.8}
Let $K$ be a local field of characteristic zero and let $\mu_n$ be the group
of $n$th roots of $1$ in $K$. Then
\[
 (K^\times:K^{\times n})=n\frac{|\mu_n|}{|n|}.
\]
If $K$ is nonarchimedean and $U$ is its unit group, then
\[
 (U:U^n)=\frac{|\mu_n|}{|n|}.
\]
\end{authorresult}
\begin{proof}
If $K=\mathbb C$, the first equation becomes $1=n\cdot n/n^2$.
If $K=\mathbb R$, it becomes $1=n\cdot1/n$ for odd $n$ and
$2=n\cdot2/n$ for even $n$.
For nonarchimedean $K$, $K^\times\simeq U\times\mathbb Z$, so
\[
 (K^\times:K^{\times n})=(U:U^n)(\mathbb Z:n\mathbb Z)=(U:U^n)n.
\]
It remains to prove the second equation. For an abelian group $M$, write
\[
 h_n(M)=\frac{(M:nM)}{(M_n:1)},\qquad M_n=\{x\in M\mid nx=0\}.
\]
This is the Herbrand quotient of $M$ regarded as a $\mathbb Z/n\mathbb Z$-module
with trivial action, so the results of II\S2 apply.
As in the proof of III.2.4, the exponential map defines an isomorphism from
a subgroup of finite index in $\mathcal O_K$ onto a subgroup of finite
index in $U$. Therefore
\[
 h_n(U)=h_n(\mathcal O_K)=(\mathcal O_K:n\mathcal O_K)
 \overset{\mathrm{def}}=|n|^{-1}.
\]
Hence $(U:U^n)=(U_n:1)/|n|$. Since $U_n=\mu_n$, the second equation follows.
\end{proof}
% END PRIMARY PROOF VII.5.1 LOCAL POWER INDEX

% BEGIN PRIMARY PROOF VII.5.1 UNIT MAP AND COMPLETION
\begin{proof}[Proof of Lemma VII.6.7]
Clearly $K^\times\cap E\supset U(S\cup T)^p$. The unit theorem gives
\[
 (U(S\cup T):U(S\cup T)^p)=p^{s+t},
\]
so it remains to prove $K^\times\cap E\subset U(S\cup T)^p$.
The next two lemmas accomplish this; the first shows that the second may
be applied to prove the inclusion.
\end{proof}

\begin{authorresult}{Lemma}{6.9}
With the above hypotheses, the obvious map
\[
 U(S)\longrightarrow\prod_{v\in T}U_v/U_v^p
\]
is surjective.
\end{authorresult}
\begin{proof}
Let $H$ be its kernel. To show surjectivity, it is enough to show
\[
 (U(S):H)=\prod_{v\in T}(U_v:U_v^p).
\]
$T$ is disjoint from $S$, so $|p|_v=1$ for $v\in T$.
Proposition VII.6.8 shows that the right side is $p^t$.
By Lemma VII.6.3, $H=U(S)\cap L^{\times p}$, whence
\[
 U(S)/H=U(S)/(U(S)\cap L^{\times p})
 \simeq U(S)L^{\times p}/L^{\times p}.
\]
By Kummer theory (A.3), the last group corresponds to $M/L$ and therefore
has order $[M:L]=p^t$.
\end{proof}

\begin{authorresult}{Proposition}{6.10}
Let $K$ be a number field containing a primitive $n$th root of $1$.
Let $S$ be a set of primes containing the infinite primes, the divisors of
$n$, and a set of representatives of the ideal class group of $K$.
Let $T$ be disjoint from $S$, and suppose that
\[
 U(S)\longrightarrow\prod_{v\in T}U_v/U_v^n
\]
is surjective. If $b\in K^\times$ is an $n$th power in $K_v$ for every
$v\in S$ and is a unit outside $S\cup T$, then $b\in K^{\times n}$.
\end{authorresult}
\begin{proof}
Let $L=K[b^{1/n}]$; we must show that $L=K$. Put
\[
 D=\prod_{v\in S}K_v^\times\times
   \prod_{v\in T}U_v^n\times\prod_{v\notin S\cup T}U_v.
\]
By Lemma VII.4.5 it suffices to prove (a) $D\subset\Nm_{L/K}\I_L$ and
(b) $DK^\times=\I_K$.

(a) Let $\mathbf d=(d_v)\in D$. Show that $d_v$ is a norm from
$K_v[b^{1/n}]$ for every $v$. If $v\in S$, then $K_v[b^{1/n}]=K_v$, so
every element of $K_v$ is a norm. If $v\in T$, local class field theory says
\[
 (K_v^\times:\Nm K_v[b^{1/n}]^\times)=[K_v[b^{1/n}]:K_v],
\]
which divides $n$, so every $n$th power in $K_v$ is a norm.
If $v\notin S\cup T$, then $nb$ is a unit at $v$; the field
$K_v[b^{1/n}]$ is therefore unramified over $K_v$, and hence every unit
is a norm.

(b) Clearly $\I_S/D=\prod_{v\in T}U_v/U_v^n$, and by hypothesis the map
$U(S)\longrightarrow\prod_{v\in T}U_v/U_v^n$ is surjective.
Thus $\I_S=D\,U(S)$ and
\[
 \I_K=\I_S K^\times=D\,U(S)K^\times=DK^\times.
\]
\end{proof}

\emph{This completes the proof of Theorem VII.5.1 (the second inequality).}
% END PRIMARY PROOF VII.5.1 UNIT MAP AND COMPLETION

\medskip
\noindent\small
The single collected target is VII.5.1. Its three conclusions, the cyclic
and Sylow reductions, the auxiliary Kummer extension and prime set, and
VII.6.1--6.10 form one proof group. The preceding first inequality and
VII.4.5--4.7 are uncounted source-local support. The source's Remark VII.5.5
and its alternative analytic route through Chapter VI are not additional
targets or additional proof routes in this reader.
\end{document}
